\section*{अध्याय \ref{ch:b1:newton-dynamics} --- गतिकी: न्यूटन के नियम}

\begin{solution}{exo:b1:newton-dynamics:1}
तराज़ू का पाठ $N = m(g + a)$, जहाँ $a$ ऊपर की ओर का त्वरण है: ऊपर,
$70 \times 11.81 = \qty{827}{N}$ (``\qty{84}{kg}''); नीचे, $70 \times 7.81 =
\qty{547}{N}$ (``\qty{56}{kg}''); अचर वेग पर \qty{687}{N}
(\qty{70}{kg}); मुक्त पतन में $0$।
\end{solution}

\begin{solution}{exo:b1:newton-dynamics:2}
दोनों के लिए एक ही त्वरण $a$ (अवितान्य डोरी): $m_1a = T$,
$m_2a = m_2g - T$, अतः $a = m_2g/(m_1 + m_2) = \qty{3.3}{m/s^2}$,
$T = m_1a = \qty{6.5}{N}$।
\end{solution}

\begin{solution}{exo:b1:newton-dynamics:3}
$a = g(\sin\alpha - \mu_d\cos\alpha) = 9.81(0.50 - 0.17) = \qty{3.2}{m/s^2}$;
$v = \sqrt{2 \times 3.2 \times 5.0} = \qty{5.7}{m/s}$। \omterm{def:b1:newton-dynamics:momentum}{द्रव्यमान} कहीं नहीं दिखता।
\end{solution}

\begin{solution}{exo:b1:newton-dynamics:4}
$k\Delta\ell = mg$: $\Delta\ell = 0.50 \times 9.81/200 = \qty{2.5}{cm}$।
साम्यावस्था से विस्थापन $x$ लेने पर $m\ddot x = -kx$ (भार और स्थैतिक खिंचाव
आपस में कट जाते हैं): $\omega_0 = \sqrt{k/m} = \qty{20}{rad/s}$,
$T = \qty{0.31}{s}$।
\end{solution}

\begin{solution}{exo:b1:newton-dynamics:5}
$\mu_s = \tan 22^\circ = 0.40$। मोड़: $v \leq \sqrt{\mu_sgR}$: $\sqrt{0.8
\times 9.81 \times 50} = \qty{20}{m/s}$ (\qty{71}{km/h}); गीली सड़क पर
\qty{14}{m/s} (\qty{50}{km/h})।
\end{solution}

\begin{solution}{exo:b1:newton-dynamics:6}
पैराशूट-यात्री: $v_\infty = \sqrt{2 \times 785/(1.2 \times 0.70)} = \qty{43}{m/s}$;
$\mathrm{Re} = \rho vL/\eta \approx 1.2 \times 43 \times 0.5/1.8\times10^{-5}
\approx 10^6$: द्विघाती। बूँद: $m = \tfrac43\pi r^3\rho_w = \qty{4.2e-12}{kg}$,
$v_\infty = mg/6\pi\eta r = 4.1\times10^{-11}/3.4\times10^{-9} = \qty{1.2}{cm/s}$;
$\mathrm{Re} = 1.2 \times 0.012 \times 2\times10^{-5}/1.8\times10^{-5} = 0.02$:
रैखिक।
\end{solution}

\begin{solution}{exo:b1:newton-dynamics:7}
तल पर: $v^2 = 2g\ell(1 - \cos 60^\circ) = g\ell$; $T = mg + mv^2/\ell =
2mg$। छोड़ने के बिंदु पर: $v = 0$, $T = mg\cos\theta_0 = 0.5mg$।
\end{solution}

\begin{solution}{exo:b1:newton-dynamics:8}
$v = \sqrt{gR\tan\beta} = \sqrt{9.81 \times 200 \times 0.268} = \qty{23}{m/s}$
(\qty{82}{km/h})। धीमे चलने पर: कार ढाल से नीचे फिसलने को होती है, घर्षण को
ढाल के ऊपर की ओर लगना पड़ता है; तेज़ चलने पर: \omterm{prop:b1:newton-dynamics:centripetal}{अभिकेंद्र बल} बढ़ाने के लिए
घर्षण ढाल के नीचे की ओर लगता है।
\end{solution}

\begin{solution}{exo:b1:newton-dynamics:9}
$T\cos\alpha = mg$, $T\sin\alpha = m\omega^2\ell\sin\alpha$: $\omega =
\sqrt{g/\ell\cos\alpha} = \sqrt{9.81/0.866} = \qty{3.4}{rad/s}$, आवर्तकाल
\qty{1.9}{s}; $T = mg/\cos\alpha = 1.15mg$।
\end{solution}

\begin{solution}{exo:b1:newton-dynamics:10}
$a = 27.8/8.0 = \qty{3.5}{m/s^2}$, $F = \qty{3.5}{kN}$; $F \leq \mu_smg$
से $\mu_s \geq a/g = 0.35$। गीली सड़क पर $\mu_s \approx 0.4$ त्वरण को
$0.4g$ के आसपास बाँध देता है, इंजन चाहे जो हो: पकड़, शक्ति नहीं।
\end{solution}

\begin{solution}{exo:b1:newton-dynamics:11}
$v = v_\infty(1 - \eu^{-t/\tau})$, $v_\infty = g\tau = \qty{4.9}{m/s}$;
$z = \int v = v_\infty[t - \tau(1 - \eu^{-t/\tau})]$। \qty{2.0}{s} पर:
$4.9 \times (2.0 - 0.5 \times 0.98) = \qty{7.4}{m}$, जबकि निर्वात में $\tfrac12 gt^2
= \qty{20}{m}$।
\end{solution}

\begin{solution}{exo:b1:newton-dynamics:12}
त्रिज्य प्रक्षेप (केंद्र की ओर): $N - mg\cos\theta = mv^2/R$, अतः
$N = m(v^2/R + g\cos\theta)$। शिखर ($\theta = \pi$): $N = m(v^2/R - g)
\geq 0$ तभी जब $v \geq \sqrt{gR}$। चाल के नियम के साथ, $v_{\text{शिखर}}^2 =
v_0^2 - 4gR \geq gR$: $v_0 \geq \sqrt{5gR}$; तब तल पर $N =
m(5g + g) = 6mg$।
\end{solution}

\begin{solution}{pb:b1:newton-dynamics:1}
\textbf{1.} $m\dot v = mg$: $v = gt$, $z = \tfrac12 gt^2$; \qty{10}{s} पर:
\qty{98}{m/s}, \qty{490}{m}।

\textbf{2.} $mg = \qty{785}{N}$।

\textbf{3.} भूमि पर फ़र्श $N = mg$ से ऊपर की ओर धकेलता है --- वही धक्का भार
के रूप में अनुभव होता है। मुक्त पतन में कोई संस्पर्श \omterm{def:b1:newton-dynamics:momentum}{बल} नहीं है: गुरुत्व
हर हिस्से पर एक-सा काम करता है और कुछ भी नहीं दबता।

\textbf{4.} एक मिनट के दौरान प्रतिशत के अंश तक हाँ; उपेक्षित: पृथ्वी का
घूर्णन (अपकेंद्र और कोरिओलिस पद,
\cref{ch:b1:non-inertial-frames}) तथा हवा।

\textbf{5.} $F_d = \tfrac12\rho C_xSv^2 = 0.48\,v^2$ (SI)।

\textbf{6.} $m\,\dd v/\dd t = mg - 0.48v^2$।

\textbf{7.} $v_\infty = \sqrt{mg/0.48} = \sqrt{1635} = \qty{40}{m/s} =
\qty{146}{km/h}$।

\textbf{8.} $mg$ से भाग दीजिए: $\dd u/\dd t\,(v_\infty/g) = 1 - u^2$;
$\tau = 40.4/9.81 = \qty{4.1}{s}$।

\textbf{9.} $\int_0^u\dd u'/(1 - u'^2) = \operatorname{artanh}u = t/\tau$,
अतः $u = \tanh(t/\tau)$।

\textbf{10.} $t = \tau\operatorname{artanh}(0.9) = 1.47\tau = \qty{6.1}{s}$;
$0.99$ के लिए: $2.65\tau = \qty{11}{s}$।

\textbf{11.} $z = \int v_\infty\tanh(t/\tau)\dd t = v_\infty\tau\ln\cosh(t/\tau)$;
\qty{10}{s} पर: $166 \times \ln\cosh(2.43) = 166 \times 1.74 = \qty{290}{m}$,
जबकि \qty{490}{m}।

\textbf{12.} $\mathrm{Re} = 1.2 \times 40 \times 0.5/1.8\times10^{-5} \approx
\num{1.3e6}$: $1000$ से कहीं ऊपर, अतः द्विघाती।

\textbf{13.} $v_\infty' = \sqrt{785/(0.6 \times 25)} = \sqrt{52} = \qty{7.2}{m/s}$।

\textbf{14.} कर्षण $= 15 \times 1600 = \qty{24}{kN}$; $a = (24000 - 785)/80
= \qty{290}{m/s^2} \approx 30g$: बंधन-पट्टा टूट जाता, या कूदने वाला।

\textbf{15.} औसत $a \approx (40 - 7)/3 = \qty{11}{m/s^2} \approx 1.1g$;
बंधन-बल $\approx m(g + a) \approx \qty{1.7}{kN}$।

\textbf{16.} $\tau' = 7.2/9.81 = \qty{0.73}{s}$: पूरा खुलने के कुछ ही
सेकंड में अवतरण स्थायी हो जाता है।

\textbf{17.} $\vect a = \vect 0$: बंधन-बल $= mg = \qty{785}{N}$।

\textbf{18.} संतुलन $mg = \tfrac12\rho C_xSv^2$ से $v \propto (C_xS)^{-1/2}$:
आधा क्षेत्रफल, $\sqrt2$ गुना चाल (\qty{10}{m/s})।

\textbf{19.} $v \propto \sqrt m$: $7.2 \times \sqrt{100/80} = \qty{8.1}{m/s}$।

\textbf{20.} $h = v^2/2g = 52/19.6 = \qty{2.6}{m}$: पहली मंज़िल से लगाई गई
छलाँग।

\textbf{21.} $a = v^2/2d$: \qty{0.5}{m} में \qty{52}{m/s^2} ($5.3g$),
\omterm{def:b1:newton-dynamics:momentum}{बल} $m(g + a) \approx \qty{5}{kN}$; \qty{1.5}{m} में \qty{17}{m/s^2} ($1.8g$),
\qty{2.2}{kN}। लुढ़ककर उतरिए।

\textbf{22.} बड़ा कर्षण (और कुछ उत्थापन) क्षण भर के लिए भार से अधिक हो
जाता है: कूदने वाला मंदित होता है और $v_\infty'$ से नीचे भूमि छूता है।

\textbf{23.} $m = \tfrac43\pi r^3\rho_w = \qty{4.2e-6}{kg}$, $S = \pi r^2 =
\qty{3.1e-6}{m^2}$: अतः $v_\infty = \sqrt{2mg/\rho C_xS} = \sqrt{8.2\times10^{-5}/
1.9\times10^{-6}} = \qty{6.6}{m/s}$;
और $\mathrm{Re} = 1.2 \times 6.6 \times 2\times10^{-3}/1.8\times10^{-5} \approx 900$:
यानी (मोटे तौर पर) द्विघाती।

\textbf{24.} $v_\infty = \qty{1.2}{cm/s}$ (\cref{exo:b1:newton-dynamics:6});
$\mathrm{Re} \approx 0.02 \ll 1$: स्टोक्स संगति के साथ लागू होता है।

\textbf{25.} $m\dot{\vect v} = m\vect g + \vect F_{\text{कर्षण}}$; कर्षण
छोटी, धीमी चीज़ों के लिए $v$ में रैखिक और बड़ी, तेज़ चीज़ों के लिए
द्विघाती। संख्याएँ: गिरते पिंड के लिए \qty{40}{m/s}, छत्र के नीचे
\qty{7}{m/s} (या वर्षा-बूँद के लिए), और कोहरे के लिए \qty{1}{cm/s} ---
वही एक नियम, परिमाण की चार कोटियों में फैला हुआ।
\end{solution}
